% Part: lambda-calculus
% Chapter: lambda-definability
% Section: introduction

\documentclass[../../../include/open-logic-section]{subfiles}

\begin{document}

\olfileid{lam}{ldf}{int}
\olsection{Pambuka}

Yen dideleng sapisan, kalkulus
lambda mung kalkulus ekspresi
kang abstrak banget kanggo
makili fungsi lan aplikasi
fungsi marang fungsi liyane.
Ora ana ing sintaksis
kalkulus lambda kang nuduhake
yen fungsi-fungsi iki gegayutan
karo jinis objek tartamtu;
mligine, sintaksise ora nyebut
bilangan asli.
Operasi dhasare---aplikasi
lan abstraksi lambda---bisa
ditrapake marang fungsi apa wae,
ora mung fungsi ing bilangan
asli.

Nanging, kanthi rekadaya,
kita bisa netepake fungsi
aritmetika, yaiku fungsi
ing bilangan asli, ing
kalkulus lambda.
Kanggo iku, saben bilangan
asli~$n \in \Nat$ diwenehi
suku-$\lambd$ mligi~$\num n$,
yaiku \emph{numeral Church}
kanggo~$n$.
(Numeral Church dijenengi
miturut Alonzo Church.)

\begin{defn}
  Yen $n \in \Nat$,
  \emph{numeral Church}
  kang cocog, $\num{n}$,
  makili~$n$:
  \[
    \num{n} \ident \lambd[fx][f^n(x)]
  \]
  Ing kene $f^n(x)$
  ateges asil ngetrapake
  $f$ marang~$x$
  ping $n$.
  Contone, $\num{0}$
  yaiku $\lambd[fx][x]$,
  lan $\num{3}$ yaiku
  $\lambd[fx][f(f(f\,x))]$.
\end{defn}
  
Numeral Church $\num n$
dikodeake minangka suku
lambda kang makili fungsi
kang nampa rong argumen,
$f$ lan $x$, banjur
ngasilake $f^n(x)$.
Numeral Church cetha
ana ing wujud normal.

Perwakilan bilangan asli
ing kalkulus lambda
mung migunani yen kita
bisa ngitung nganggo
bilangan kasebut.
Ngitung nganggo numeral
Church ing kalkulus lambda
ateges ngetrapake
suku-$\lambd$~$F$ marang
numeral Church mau,
banjur ngreduksi suku
gabungan~$F\, \num n$
nganti wujud normal.
Yen suku iki mesthi
direduksi nganti wujud
normal lan wujud normale
mesthi numeral Church~$\num m$,
asil pangétungan bisa
dianggep bilangan~$m$.
Mula $F$ bisa dianggep
netepake fungsi
$f\colon \Nat \to \Nat$,
yaiku fungsi kang
$f(n) = m$ yen lan
mung yen $F\, \num n
\red \num m$.
Miturut sipat Church--Rosser,
wujud normal iku tunggal
yen ana. Dadi, yen
$F\, \num n \red \num m$,
ora ana suku wujud normal
liya, mligine ora ana
numeral Church liya,
kang bisa digayuh saka
$F \, \num n$ kanthi
reduksi.

Kosok baline, kanggo
fungsi $f\colon \Nat \to \Nat$,
kita bisa takon apa ana
suku $F$ kang netepake~$f$
kanthi cara iki.
Yen ana, kita kandha
yen $F$ \emph{!!{lambda define}s}~$f$,
lan $f$ iku
!!{lambda definable}.
Konsep iki bisa ditambahi
marang fungsi kang aritase
luwih saka siji lan fungsi
parsial.

\begin{defn}
Upamana $f\colon \Nat^k \to \Nat$.
Suku lambda~$F$ diarani
\emph{!!{lambda define}s}
$f$ yen kanggo saben
$n_0$, \dots,~$n_{k-1}$,
\[
F \, \num{n_0} \, \num{n_1} \dots \num{n_{k-1}} \red \num{f(n_0, n_1, \dots,
  n_{k-1})}
\]
yen $f(n_0, \dots, n_{k-1})$
katetepake, lan
$F \, \num{n_0} \,
\num{n_1} \dots \num{n_{k-1}}$
ora nduweni wujud normal
yen fungsine ora katetepake.
\end{defn}

Conto kang prasaja banget
yaiku fungsi konstan.
Suku $C_k \ident
\lambd[x][\num{k}]$
!!{lambda define}s
fungsi $c_k\colon \Nat \to
\Nat$ kang nyukupi
$c_k(n) = k$.
Jalaran, $C_k \, \num n \ident
(\lambd[x][\num{k}])\num n
\redone \num{k}$
kanggo sembarang~$n$.
Fungsi identitas
!!{lambda defined}
dening $\lambd[x][x]$.
Fungsi kang luwih rumit
mesthi luwih angel
ditetepake lan asring
mbutuhake rekadaya akeh.
Mula rada nggumunake yen
saben fungsi kang bisa
diitung iku
!!{lambda definable}.
Kosok baline uga bener:
yen fungsi iku
!!{lambda definable},
fungsi iku bisa diitung.

\end{document}
